% GENERATED FILE. Edit canonical lessons, then run npm run generate:books.
% canonical-source-hash: fnv1a64:9509407cc0308b51
% canonical-lessons: ZH-C263-jingli, ZH-C263-fuze, ZH-C263-qingjia, ZH-W263-jia-leave, ZH-C263-bingjia, ZH-C263-nianjia

\chapter{Time Off Work}
\label{ch:zh-time-off-work}
\hlchaptermodality{263}

\begin{chapteropening}
\textbf{By the end of this chapter:} \emph{I can talk about the manager, say who is in charge, and ask for time off, sick leave or annual leave.}

Complete the last lesson of chapter 263: I can talk about the manager, say who is in charge, and ask for time off, sick leave or annual leave.
\end{chapteropening}

\section[jīnglǐ]{jīnglǐ — manager}

\label{lesson:ZH-C263-jingli}



\begin{quote}
Before the new one: say the Mandarin for to defeat, then the Mandarin for to cheer someone on.
\end{quote}

\subsection*{You'll want to know: jīnglǐ}
\emph{jīnglǐ} — \textquotedblleft{}manager\textquotedblright{}.

\begin{quote}
\emph{jīnglǐ jīntiān bú zài} — The manager isn't in today.
\end{quote}

Its characters come later; for now it is a word for the ear.

\begin{cousinweb}
\emph{jīnglǐ} is built from \emph{jīng} and \emph{lǐ}, both inherited from Middle Chinese, \emph{keng} and \emph{liX} (Guangyun, in Baxter's transcription), which glosses the first 'constant' and the second 'to put in order', and quotes an older dictionary for the second, 'to work jade'.
\end{cousinweb}

\begin{grammarlens}[title={What you've built}]
One more word for saying what I do.
\end{grammarlens}

\subsection*{Your turn}
\begin{itemize}
\raggedright
  \item \emph{Say it:} \emph{jīnglǐ}
  \item \emph{Say it:} \emph{jīnglǐ}, once more
  \item \emph{Say it:} \emph{jīnglǐ jīntiān bú zài} — The manager isn't in today.
  \item \emph{Recall:} say the Mandarin for to defeat, then the Mandarin for to cheer someone on, then say \emph{jīnglǐ} again
\end{itemize}

\begin{tcolorbox}[breakable,colback=teal!4,colframe=teal!35!black,title={Before you move on}]
What does jīnglǐ mean? (\textquotedblleft{}Manager\textquotedblright{}.) Say it once more.
\end{tcolorbox}
\section[fùzé]{fùzé — to be in charge (of)}

\label{lesson:ZH-C263-fuze}



\begin{quote}
Before the new one: say the Mandarin for to cheer someone on, then the Mandarin for a manager.
\end{quote}

\subsection*{You'll want to know: fùzé}
\emph{fùzé} — \textquotedblleft{}to be in charge (of)\textquotedblright{}.

\begin{quote}
\emph{zhège gōngzuò wǒ fùzé} — I'm in charge of this job.
\end{quote}

Its characters come later; for now it is a word for the ear.

\begin{cousinweb}
\emph{fùzé} is built from \emph{fù} 'to carry' and \emph{zé} 'to demand', both inherited from Middle Chinese, \emph{bjuwX} and \emph{tsreak} (Guangyun, in Baxter's transcription), which glosses the first 'to bear on the back' and the second 'to seek'.
\end{cousinweb}

\begin{grammarlens}[title={What you've built}]
One more word for saying what I do.
\end{grammarlens}

\subsection*{Your turn}
\begin{itemize}
\raggedright
  \item \emph{Say it:} \emph{fùzé}
  \item \emph{Say it:} \emph{fùzé}, once more
  \item \emph{Say it:} \emph{zhège gōngzuò wǒ fùzé} — I'm in charge of this job.
  \item \emph{Recall:} say the Mandarin for to cheer someone on, then the Mandarin for a manager, then say \emph{fùzé} again
\end{itemize}

\begin{tcolorbox}[breakable,colback=teal!4,colframe=teal!35!black,title={Before you move on}]
What does fùzé mean? (\textquotedblleft{}To be in charge (of)\textquotedblright{}.) Say it once more.
\end{tcolorbox}
\section[\texorpdfstring{\zh{请假} (qǐngjià)}{请假 (qǐngjià)}]{\zh{请假} (qǐngjià) — to ask for time off}

\label{lesson:ZH-C263-qingjia}



\begin{quote}
Before the new one: say the Mandarin for a manager, then the Mandarin for to be in charge.
\end{quote}

\subsection*{You'll want to know: \zh{请假}}
\textbf{\zh{请假}} — \emph{qǐngjià} — \textquotedblleft{}to ask for time off\textquotedblright{}.

\begin{quote}
\textbf{\zh{我明天要请假}} \emph{wǒ míngtiān yào qǐngjià} — I need to take tomorrow off.
\end{quote}

The new character, in pieces: \zh{假} — eleven strokes: the person piece \zh{亻}; then a narrow piece of five strokes, a short bar that turns down, a short bar, a long vertical and two short bars against it; then four strokes, a short bar that turns down, a short bar, a stroke that runs right and sweeps down to the left, and a long stroke to the right across it. The next lesson puts it in your hand.

\begin{cousinweb}
\textbf{\zh{请假}} (\emph{qǐngjià}) is built from \zh{请} 'to ask' and \zh{假} 'leave', both inherited from Middle Chinese, \emph{tshjengX} and \emph{kaeH} (Guangyun, in Baxter's transcription), which glosses the first 'to beg' and the second 'leave from work'.
\end{cousinweb}

\begin{grammarlens}[title={What you've built}]
One more word for saying what I do.
\end{grammarlens}

\subsection*{Your turn}
\begin{itemize}
\raggedright
  \item \emph{Say it:} \emph{qǐngjià}
  \item \emph{Say it:} \emph{qǐngjià}, once more
  \item \emph{Say it:} \emph{wǒ míngtiān yào qǐngjià} — I need to take tomorrow off.
  \item \emph{Recall:} say the Mandarin for a manager, then the Mandarin for to be in charge, then say \emph{qǐngjià} again
\end{itemize}

\begin{tcolorbox}[breakable,colback=teal!4,colframe=teal!35!black,title={Before you move on}]
What does \zh{请假} mean? (\textquotedblleft{}To ask for time off\textquotedblright{}.) Say it once more.
\end{tcolorbox}
\section[\texorpdfstring{\zh{假} (jià)}{假 (jià)}]{\zh{假} — the person piece, then two narrow pieces}

\label{lesson:ZH-W263-jia-leave}



\begin{quote}
The word first, then the character inside it.

\begin{itemize}
\raggedright
  \item \emph{Recall:} say \emph{qǐngjià} — \textbf{R1}, one lesson back, before it decays
\end{itemize}

\begin{quote}
\zh{请} \zh{你}
\end{quote}

\emph{Read it:} each one aloud, with nothing said first — \emph{qǐng}, \emph{nǐ}
\end{quote}

\subsection*{Script — the person piece, then two narrow pieces}
\hlblockfigure{\detokenize{figures/ZH-W263-jia-leave-filmstrip.pdf}}{How \zh{假} is written, stroke by stroke}

\begin{quote}
\zh{假}
\end{quote}

\textbf{Three pieces, left to right.}

\textbf{The left piece, \zh{亻}, two strokes:} a stroke falling to the left, then a vertical.

\textbf{The middle piece, five strokes.} A short bar that turns down at its right end. A short bar under it. A long vertical down the left side to the foot. Then two short bars against the vertical, the upper one first.

\textbf{The right piece, four strokes.} A short bar that turns down. A short bar under it. A stroke that runs right and sweeps down to the left. Then a long stroke falling to the right across it to the foot.

2 + 5 + 4 = eleven.

\subsection*{Your turn}
\begin{itemize}
\raggedright
  \item \emph{Write it:} \zh{假}, all eleven, with the model in front of you
  \item \emph{Cover:} write it again from memory, and count the pieces, not the strokes
\end{itemize}

\begin{tcolorbox}[breakable,colback=teal!4,colframe=teal!35!black,title={Before you move on}]
How many strokes in \zh{假}? (\textbf{Eleven}.) Which piece comes first? (\textbf{The person piece, \zh{亻}}.) What is the last stroke? (\textbf{A long stroke falling to the right}.)

Source: \href{https://raw.githubusercontent.com/chanind/hanzi-writer-data/68d10a4b21150cae5e1ebbd223eed289cf32d90c/data/\%E5\%81\%87.json}{Hanzi Writer Data \zh{假}, pinned stroke paths and medians}.
\end{tcolorbox}
\section[\texorpdfstring{\zh{病假} (bìngjià)}{病假 (bìngjià)}]{\zh{病假} (bìngjià) — sick leave}

\label{lesson:ZH-C263-bingjia}



\begin{quote}
Before the new one: say the Mandarin for to be in charge, then the Mandarin for to ask for time off.
\end{quote}

\subsection*{You'll want to know: \zh{病假}}
\textbf{\zh{病假}} — \emph{bìngjià} — \textquotedblleft{}sick leave\textquotedblright{}.

\begin{quote}
\textbf{\zh{他今天请了病假}} \emph{tā jīntiān qǐngle bìngjià} — He's off sick today.
\end{quote}

\begin{cousinweb}
\textbf{\zh{病假}} (\emph{bìngjià}) is built from \zh{病} 'illness' and \zh{假} 'leave', both inherited from Middle Chinese, \emph{bjaengH} and \emph{kaeH} (Guangyun, in Baxter's transcription), which glosses the first 'grief' and quotes an older dictionary for 'an illness getting worse'.
\end{cousinweb}

\begin{grammarlens}[title={What you've built}]
One more word for saying what I do.
\end{grammarlens}

\subsection*{Your turn}
\begin{itemize}
\raggedright
  \item \emph{Say it:} \emph{bìngjià}
  \item \emph{Say it:} \emph{bìngjià}, once more
  \item \emph{Say it:} \emph{tā jīntiān qǐngle bìngjià} — He's off sick today.
  \item \emph{Recall:} say the Mandarin for to be in charge, then the Mandarin for to ask for time off, then say \emph{bìngjià} again
\end{itemize}

\begin{tcolorbox}[breakable,colback=teal!4,colframe=teal!35!black,title={Before you move on}]
What does \zh{病假} mean? (\textquotedblleft{}Sick leave\textquotedblright{}.) Say it once more.
\end{tcolorbox}
\section[\texorpdfstring{\zh{年假} (niánjià)}{年假 (niánjià)}]{\zh{年假} (niánjià) — annual leave}

\label{lesson:ZH-C263-nianjia}



\begin{quote}
Before the new one: say the Mandarin for to ask for time off, then the Mandarin for sick leave.
\end{quote}

\subsection*{You'll want to know: \zh{年假}}
\textbf{\zh{年假}} — \emph{niánjià} — \textquotedblleft{}annual leave\textquotedblright{}.

\begin{quote}
\textbf{\zh{我今年的年假有十天}} \emph{wǒ jīnnián de niánjià yǒu shí tiān} — I get ten days' leave this year.
\end{quote}

\begin{cousinweb}
\textbf{\zh{年假}} (\emph{niánjià}) is built from \zh{年} 'a year' and \zh{假} 'leave', both inherited from Middle Chinese, \emph{nen} and \emph{kaeH} (Guangyun, in Baxter's transcription).
\end{cousinweb}

\begin{grammarlens}[title={What you've built}]
One more word for saying what I do.
\end{grammarlens}

\subsection*{Your turn}
\begin{itemize}
\raggedright
  \item \emph{Say it:} \emph{niánjià}
  \item \emph{Say it:} \emph{niánjià}, once more
  \item \emph{Say it:} \emph{wǒ jīnnián de niánjià yǒu shí tiān} — I get ten days' leave this year.
  \item \emph{Recall:} say the Mandarin for to ask for time off, then the Mandarin for sick leave, then say \emph{niánjià} again
\end{itemize}

\begin{tcolorbox}[breakable,colback=teal!4,colframe=teal!35!black,title={Before you move on}]
What does \zh{年假} mean? (\textquotedblleft{}Annual leave\textquotedblright{}.) Say it once more.
\end{tcolorbox}
